\documentclass[fontset=mac]{ctexart}

% ====== 页面与数学宏包 ======
\usepackage[top=2.54cm, bottom=2.54cm, left=3.18cm, right=3.18cm]{geometry}
\usepackage{amsmath, amssymb, amsthm}
\usepackage{bm}
\usepackage{hyperref}
\usepackage{booktabs}
\usepackage{graphicx}
\usepackage{float}
\usepackage{array}

% ====== 量子算符 ======
\newcommand{\adag}[1]{\hat{a}_{#1}^\dagger}
\newcommand{\aop}[1]{\hat{a}_{#1}}
\newcommand{\Xop}{X}
\newcommand{\Yop}{Y}
\newcommand{\Zop}{Z}
\newcommand{\Iop}{I}
\newcommand{\Nop}[1]{\hat{N}_{#1}}
\newcommand{\ket}[1]{|#1\rangle}
\newcommand{\bra}[1]{\langle #1|}
\newcommand{\braket}[2]{\langle #1|#2\rangle}
\newcommand{\HS}{\mathcal{H}}
\newcommand{\Hred}{\mathcal{H}_{\text{red}}}

% ====== 标题 ======
\title{Problem 2.2: HF State and SQD Subspace Diagonalization}
\author{nyloong}
\date{\today}

\begin{document}
\maketitle

% ===========================================================================
\section*{Problem Statement}

SQD 从一个量子态 $\ket{\Psi}$ 中采样 $S$ 个比特串 (bitstrings)，每个比特串对应一个电子组态
(Slater 行列式 $\ket{D_k}$)，然后在子空间 $\mathcal{S}= \operatorname{span}\{\ket{D_k}\}$ 中精确对角化 CI 矩阵
$H_{kl} = \braket{D_k}{\hat{H}}{D_l}$。最简单的初始态是 HF 态
$\ket{\text{HF}} = \prod_{i\in\text{occ}} \adag{i} \ket{0}$——它是一个单一的 computational basis state。

\textbf{给定分子:} (a) H$_2$, 4 个自旋轨道 (指标 $0$--$3$), 2 个电子。

\textbf{求解:}
\begin{enumerate}
    \item[(a)] 确定 H$_2$ 的 HF 态在 JW 映射下对应哪个 computational basis state $\ket{b_3 b_2 b_1 b_0}$ ($b_i\in\{0,1\}$),
    并写出其制备电路。手动列出所有 $\binom{4}{2} = 6$ 个组态，构造 $6\times 6$ CI 矩阵，精确对角化，
    与 PySCF FCI 能量比较。
    \item[(b)] 若子空间仅包含 HF + 所有单激发 (无双激发)，SQD 能量等于什么？用 Slater-Condon 规则
    和 HF 自洽性解释原因。
\end{enumerate}

\textbf{Advanced Challenge:} 将 H$_2$ 键长拉伸至 3.0\,\AA。
\begin{itemize}
    \item[A.] HF + singles 的结论是否仍然成立？HF 自洽性在强关联区域会怎样？
\end{itemize}

\newpage

% ===========================================================================
\section*{(a) H$_2$/STO-3G 的 JW 映射与 CI 对角化}

\subsection*{HF 态在 JW 映射下的表示}

H$_2$/STO-3G 有 2 个空间轨道 ($\sigma_g$ 成键 + $\sigma_u^*$ 反键)，共 4 个自旋轨道。
在当前 convention 下，自旋轨道指标与 qubit 的对应关系为：

\begin{itemize}
    \item qubit 0: $\sigma_u^*$ $\beta$ 自旋 (虚轨道)
    \item qubit 1: $\sigma_u^*$ $\alpha$ 自旋 (虚轨道)
    \item qubit 2: $\sigma_g$ $\beta$ 自旋 (占据轨道)
    \item qubit 3: $\sigma_g$ $\alpha$ 自旋 (占据轨道)
\end{itemize}

RHF 基态将 2 个电子分别填入能量最低的 $\alpha$ 和 $\beta$ 自旋轨道 (两个 $\sigma_g$)，
因此在 JW 映射下对应于 qubit 2 和 3 被占据：

\begin{equation}
    \boxed{\ket{\text{HF}} = \ket{1_3 1_2 0_1 0_0} = \ket{1100}}
\end{equation}

\subsection*{制备电路}

HF 态是最简单的 computational basis state——只需在 occupied qubit 上施加 $\Xop$ 门：

\begin{equation}
    \ket{1100} = \Xop_3 \Xop_2 \ket{0000}
\end{equation}

这是一个深度为 1 的单 qubit 门电路，不需要任何纠缠门。

\subsection*{6 个二电子组态}

H$_2$ 共有 4 个自旋轨道、2 个电子，所有可能的电子组态 ($N_e=2$) 共 $\binom{4}{2}=6$ 个：

\begin{table}[H]
\centering
\begin{tabular}{c|c|c|c}
\toprule
$\ket{D_k}$ & 整数 & 占据轨道 & 激发类型 \\
\midrule
$\ket{0011}$ &  3 & $\{0,1\}$ & 双激发 (两个都到 $\sigma_u^*$) \\
$\ket{0101}$ &  5 & $\{0,2\}$ & 单激发 ($\beta$: $\sigma_g\to\sigma_u^*$) \\
$\ket{0110}$ &  6 & $\{1,2\}$ & 单激发 ($\alpha$: $\sigma_g\to\sigma_u^*$) \\
$\ket{1001}$ &  9 & $\{0,3\}$ & 单激发 ($\beta$: $\sigma_g\to\sigma_u^*$) \\
$\ket{1010}$ & 10 & $\{1,3\}$ & 单激发 ($\alpha$: $\sigma_g\to\sigma_u^*$) \\
$\ket{1100}$ & 12 & $\{2,3\}$ & HF (基态) \\
\bottomrule
\end{tabular}
\caption{H$_2$/STO-3G 的所有 6 个二电子组态}
\end{table}

注意：在 qubit 表示 $\ket{b_3 b_2 b_1 b_0}$ 中，$b_i=1$ 表示自旋轨道 $i$ 被占据。

\subsection*{$6\times 6$ CI 矩阵与精确对角化}

利用 OpenFermion + PySCF 在 RHF 轨道基下构造哈密顿量的 JW 映射，
提取二电子子空间的 $6\times 6$ CI 矩阵，精确对角化。
计算结果如下：

% --- 平衡键长 (0.74 Å) ---
\textbf{平衡键长 $R=0.74$\,\AA}：

RHF 计算结果：
\begin{align*}
    E_{\text{RHF}} &= -1.1167593074\;\text{Ha}, \\
    \varepsilon_0(\sigma_g) &= -0.57855386\;\text{Ha}, \\
    \varepsilon_1(\sigma_u^*) &= +0.67114349\;\text{Ha}, \\
    \Delta\varepsilon &= 1.2497\;\text{Ha}.
\end{align*}

$6\times 6$ CI 矩阵的形式 (以 $\ket{D_0},\ldots,\ket{D_5}$ 为基, 单位 Ha):

\begin{equation}
H_{\text{CI}} = \begin{pmatrix}
-0.2706 &  0      &  0      &  0      &  0      & -0.1812 \\
 0      & -0.4764 &  0      &  0      &  0      &  0      \\
 0      &  0      & -0.4764 &  0      &  0      &  0      \\
 0      &  0      &  0      & -0.4764 &  0      &  0      \\
 0      &  0      &  0      &  0      & -0.4764 &  0      \\
-0.1812 &  0      &  0      &  0      &  0      & -1.1168
\end{pmatrix}
\end{equation}

注意 $H_{\text{CI}}$ 的两个关键特征：
\begin{itemize}
    \item \textbf{HF--单激发块为零:} $H_{15}=H_{25}=H_{35}=H_{45}=0$。HF 态 ($\ket{D_5}$) 不与任何单激发行列式耦合。
    \item \textbf{HF--双激发耦合非零:} $H_{05}=H_{50} = -0.1812$\,\text{Ha}$\neq 0$。这正是 $\braket{\sigma_g\sigma_g}{\sigma_u^*\sigma_u^*}$ 双电子积分。
\end{itemize}

精确对角化得到 6 个本征值，最低本征值即 SQD 基态能量：

\begin{equation}
    \boxed{E_{\text{SQD}} = -1.1372838345\;\text{Ha} = E_{\text{FCI}}}
\end{equation}

与 PySCF FCI 的差异 $\Delta = 1.11\times 10^{-16}$\,\text{Ha} (机器精度)。
SQD 在包含所有二电子组态的子空间中对角化，等价于 FCI。

\subsection*{基态波函数组成}

\begin{align}
    \ket{\Psi_0} &= -0.993647\;\ket{1100}_{\text{HF}} + 0.112544\;\ket{0011}_{\text{double}} \\
    |c_{\text{HF}}|^2 &= 98.73\%,\quad |c_{\text{double}}|^2 = 1.27\%
\end{align}

所有单激发组态的系数严格为零 (对称性 + Brillouin 定理)。HF 态几乎完全主导基态波函数。

\newpage

% ===========================================================================
\section*{(b) HF + Singles 子空间与 Brillouin 定理}

\subsection*{$5\times 5$ 子空间对角化}

若子空间仅包含 HF + 4 个单激发 ($\binom{4}{2}=6$ 中去掉唯一的双激发 $\ket{0011}$)：

\begin{table}[H]
\centering
\begin{tabular}{c|c}
\toprule
态 & 含义 \\
\midrule
$\ket{D_5} = \ket{1100}$ & HF 基态 \\
$\ket{D_2} = \ket{0110}$ & $\alpha$ 单激发: $\sigma_g^\alpha \to \sigma_u^{*\alpha}$ \\
$\ket{D_4} = \ket{1010}$ & $\alpha$ 单激发: $\sigma_g^\alpha \to \sigma_u^{*\alpha}$ \\
$\ket{D_1} = \ket{0101}$ & $\beta$ 单激发: $\sigma_g^\beta \to \sigma_u^{*\beta}$ \\
$\ket{D_3} = \ket{1001}$ & $\beta$ 单激发: $\sigma_g^\beta \to \sigma_u^{*\beta}$ \\
\bottomrule
\end{tabular}
\end{table}

HF + singles 子矩阵是严格块对角的：

\begin{equation}
H_{\text{HF+singles}} = \begin{pmatrix}
E_{\text{RHF}} & \mathbf{0}^\top \\
\mathbf{0} & H_{\text{singles}}
\end{pmatrix}
= \begin{pmatrix}
-1.116759 & 0 & 0 & 0 & 0 \\
0 & -0.4764 & 0 & 0 & 0 \\
0 & 0 & -0.4764 & 0 & 0 \\
0 & 0 & 0 & -0.4764 & 0 \\
0 & 0 & 0 & 0 & -0.4764
\end{pmatrix}
\end{equation}

因此：

\begin{equation}
    \boxed{E(\text{HF+singles}) = E_{\text{RHF}} = -1.1167593074\;\text{Ha}}
\end{equation}

若 SQD 只采样到 HF + 单激发组态，基态能量就是 HF 能量——单激发不提供任何关联能修正。

\subsection*{用 Slater-Condon 规则解释}

根据 Slater-Condon 规则，HF 基态 $\ket{\Phi_0}$ 和单激发行列式 $\ket{\Phi_i^a}$ (占据轨道 $i\to$ 虚轨道 $a$) 之间的哈密顿矩阵元为：

\begin{equation}
\braket{\Phi_0}{\hat{H}}{\Phi_i^a}
= \braket{i}{\hat{h}}{a} + \sum_{j}^{\text{occ}} \Bigl[\braket{ij}{aj} - \braket{ij}{ja}\Bigr]
\end{equation}

等式右边恰好是 Fock 矩阵的占据--虚块元 $f_{ia}$。HF 方程 $\hat{f}\ket{\psi_p} = \varepsilon_p\ket{\psi_p}$ 在自洽收敛后给出 Fock 矩阵对角化：

\begin{equation}
f_{pq} = \varepsilon_p \delta_{pq}
\end{equation}

特别地，对于占据轨道 $i$ 和虚轨道 $a$：$f_{ia}=0$。

因此：

\begin{equation}
\boxed{\braket{\Phi_0}{\hat{H}}{\Phi_i^a} = f_{ia} = 0}
\end{equation}

这就是 \textbf{Brillouin 定理}：HF 基态与所有单激发组态之间的哈密顿矩阵元恒为零。
这不是一个近似，而是 HF 自洽场收敛条件的直接数学推论。

\textbf{物理后果：} 在 HF + singles 子空间中对角化，HF 态仍是该子空间的一个本征态，其本征值 $E_{\text{RHF}}$ 不变。单激发不能改善基态能量——关联能的修正必须从双激发开始。
事实上，最低阶的关联能修正来自双激发组态与 HF 的耦合：

\begin{equation}
\braket{\Phi_0}{\hat{H}}{\Phi_{ij}^{ab}} = \braket{ij}{ab} - \braket{ij}{ba} = \braket{ij\|ab} \neq 0
\end{equation}

这正是 MP2 微扰理论和耦合簇方法的基础。

\newpage

% ===========================================================================
\section*{Advanced Challenge: H$_2$ 键长拉伸至 3.0\,\AA}

\subsection*{A. HF + Singles 的结论是否仍然成立？}

\textbf{是的，Brillouin 定理在任何键长下都成立}，因为它是 HF 方程 $f_{ia}=0$ 的数学推论，
而非物理近似。

在 $R=3.0$\,\AA\, 处重复相同计算：

\begin{align*}
    E_{\text{RHF}} &= -0.6560482511\;\text{Ha}, \\
    \varepsilon_0(\sigma_g) &= -0.18053922\;\text{Ha}, \\
    \varepsilon_1(\sigma_u^*) &= +0.01807133\;\text{Ha}, \\
    \Delta\varepsilon &= 0.1986\;\text{Ha} \quad (\text{平衡键长时 } 1.2497\;\text{Ha}).
\end{align*}

$6\times 6$ CI 矩阵：

\begin{equation}
H_{\text{CI}}^{(3.0\,\text{\AA})} = \begin{pmatrix}
-0.0478 &  0      &  0      &  0      &  0      & -0.1812 \\
 0      & -0.3524 &  0      &  0      &  0      &  0      \\
 0      &  0      & -0.3524 &  0      &  0      &  0      \\
 0      &  0      &  0      & -0.3524 &  0      &  0      \\
 0      &  0      &  0      &  0      & -0.3524 &  0      \\
-0.1812 &  0      &  0      &  0      &  0      & -0.6560
\end{pmatrix}
\end{equation}

\textbf{验证 Brillouin 定理仍成立：}
\begin{equation}
    \max_{i\in\text{occ},\,a\in\text{virt}} |\braket{\text{HF}}{\hat{H}}{\text{single}}|
    < 10^{-15}\;\text{Ha} \quad \checkmark
\end{equation}

HF--单激发矩阵元仍然严格为零，$5\times 5$ 矩阵仍然是块对角的。

\textbf{但一个重要的区别出现了：} 在 $3.0$\,\AA\, 处，singles 块的本征值 ($-0.3524$\,\text{Ha}) 已经低于 $E_{\text{RHF}}$ ($-0.6560$\,\text{Ha} > $-0.3524$\,\text{Ha})！
因此在 $5\times 5$ 子空间中对角化，最低本征值不再等于 $E_{\text{RHF}}$：

\begin{equation}
    E(\text{HF+singles}) = -0.932936\;\text{Ha} < E_{\text{RHF}} = -0.656048\;\text{Ha}
\end{equation}

但这\textbf{不违反 Brillouin 定理}。定理只保证 HF 不与 singles 耦合
($\braket{\text{HF}}{\hat{H}}{\text{single}}=0$)，但 singles 块内部可以有低于 $E_{\text{RHF}}$ 的本征值。
这反映了自旋翻转激发 (spin-flip excitations) 在近简并轨道区域变得非常低能。

不过，即使 singles 块能量更低，关联能仍无法恢复：
$6\times 6$ CI 基态能量是 $-0.933632$\,\text{Ha}，
减去 $E(\text{HF+singles})$ 之差约 $0.0007$\,\text{Ha}
来自双激发组态。

\subsection*{HF 自洽性在强关联区域的失效}

强关联的本质不在于电子之间的瞬时 Coulomb 排斥 (动态关联)，而在于\textbf{轨道近简并}导致的波函数多参考特性。
在 $R=3.0$\,\AA\, 处，$\sigma_g$ 和 $\sigma_u^*$ 的能隙仅余 $0.20$\,\text{Ha}（平衡键长为 $1.25$\,\text{Ha}），
两个组态 $\ket{\sigma_g\bar{\sigma}_g}$ 和 $\ket{\sigma_u\bar{\sigma}_u}$ 近似简并。

此时基态波函数是多参考的：

\begin{align}
    \ket{\Psi_0^{(3.0\,\text{\AA})}}
    &= -0.733106\;\ket{1100}_{\text{HF}} + 0.680115\;\ket{0011}_{\text{double}} \\
    |c_{\text{HF}}|^2 &= 53.74\%,\quad |c_{\text{double}}|^2 = 46.26\%
\end{align}

HF 态仅占 $\sim 54\%$ 权重——单参考方法的出发点已经\textbf{定性错误}。

\begin{table}[H]
\centering
\begin{tabular}{lcc}
\toprule
\textbf{指标} & \textbf{平衡键长 (0.74\,\AA)} & \textbf{拉伸 (3.0\,\AA)} \\
\midrule
$\Delta\varepsilon$ (HOMO--LUMO gap) & 1.2497 Ha & 0.1986 Ha \\
$E_{\text{RHF}}$ & $-1.116759$ Ha & $-0.656048$ Ha \\
$E_{\text{FCI}}$ & $-1.137284$ Ha & $-0.933632$ Ha \\
RHF 误差 & 0.559 eV & 7.553 eV \\
HF 权重 $|c_{\text{HF}}|^2$ & 98.73\% & 53.74\% \\
双激发权重 $|c_{\text{double}}|^2$ & 1.27\% & 46.26\% \\
$\max|\braket{\text{HF}}{\hat{H}}{\text{single}}|$ & $<10^{-15}$ $\checkmark$ & $<10^{-15}$ $\checkmark$ \\
$E(\text{HF+singles}) = E_{\text{RHF}}$? & $\checkmark$ (是) & $\times$ (否, singles 块更低) \\
\bottomrule
\end{tabular}
\caption{平衡键长与拉伸键长的全面对比}
\end{table}

\subsection*{RHF 定性失败的原因}

RHF 波函数强制两个电子占据同一空间轨道 $\sigma_g$。在解离极限下：

\begin{align}
    \ket{\text{RHF}} &= \frac{1}{2}\Bigl[
        \underbrace{\ket{1s_A(1)\;1s_A(2)}}_{\text{H}^- \text{H}^+}
      + \underbrace{\ket{1s_B(1)\;1s_B(2)}}_{\text{H}^+ \text{H}^-}
      + \underbrace{\ket{1s_A(1)\;1s_B(2)} + \ket{1s_B(1)\;1s_A(2)}}_{\text{H}\cdot\;\text{H}\cdot}
    \Bigr] \\
    &= 50\%\;\text{离子项}\;+\;50\%\;\text{共价项}
\end{align}

RHF 预测解离产物为一半离子 (H$^-$H$^+$) 一半共价 (H$\cdot$H$\cdot$)，这在化学上是\textbf{荒谬的}——真实解离应为 100\% 中性氢原子。

\subsection*{UHF: 通过破坏自旋对称性获取正确物理}

UHF 允许 $\alpha$ 和 $\beta$ 电子使用\textbf{不同的空间轨道}，通过 HOMO--LUMO 混合实现空间分离：

\begin{align}
    \psi^\alpha &\approx c_A^\alpha\,\phi_{1s}^A + c_B^\alpha\,\phi_{1s}^B, \quad c_A^\alpha > c_B^\alpha \quad (\alpha\text{ 偏向 A}) \\
    \psi^\beta  &\approx c_A^\beta\,\phi_{1s}^A + c_B^\beta\,\phi_{1s}^B, \quad c_A^\beta < c_B^\beta \quad (\beta\text{ 偏向 B})
\end{align}

需用破缺对称的初始猜测 (HOMO--LUMO 混合相反符号) 启动 SCF 迭代：

\begin{table}[H]
\centering
\begin{tabular}{lccc}
\toprule
\textbf{方法} & \textbf{能量 (Ha)} & \textbf{误差 (eV)} & \textbf{$\langle\hat{S}^2\rangle$} \\
\midrule
RHF & $-0.656048$ & 7.553 & 0.0000 \\
UHF & $-0.933285$ & 0.009 & 0.9986 \\
FCI & $-0.933632$ & -- & 0.0000 \\
\bottomrule
\end{tabular}
\caption{RHF vs UHF vs FCI 在 $R=3.0$\,\AA\, 处的对比}
\end{table}

UHF 通过破坏自旋对称性，几乎完全恢复了正确的解离能量 (误差仅 $0.009$\,\text{eV})。
但代价是自旋污染：标称单重态的 $\langle\hat{S}^2\rangle \approx 1.0$，说明波函数混入了约 50\% 的三重态成分。

\textbf{多参考方法的必要性：} 真正的解决方案是使用 CASCI/CASSCF 等多参考方法，
将 $\ket{\sigma_g\bar{\sigma}_g}$ 和 $\ket{\sigma_u\bar{\sigma}_u}$ 作为\textbf{参考空间}，
在此之上叠加动力学关联修正 (如 CASPT2, NEVPT2)。这样才能同时获得正确的能量和正确的自旋本征态。

\subsection*{总结：Brillouin 定理的普适性与局限性}

Brillouin 定理 ($\braket{\Phi_0}{\hat{H}}{\Phi_i^a}=f_{ia}=0$) 是 HF 自洽条件的数学推论，
在任何键长下都严格成立。但它只保证单激发不耦合到 HF 态，并不保证 HF 本身是好的近似。

\begin{center}
\begin{tabular}{lcc}
\toprule
& \textbf{弱关联 (0.74\,\AA)} & \textbf{强关联 (3.0\,\AA)} \\
\midrule
Brillouin 定理 & $\checkmark$ 成立 & $\checkmark$ 成立 \\
HF 近似质量 & 优秀 (98.7\% 权重) & 失败 (53.7\% 权重) \\
单激发作用 & 无 (被 Brillouin 消耦合) & 无 (同左) \\
双激发作用 & 小修正 (1.3\%) & \textbf{必需} (46.3\%) \\
所需方法 & MP2/CCSD & CASCI/CASSCF \\
\bottomrule
\end{tabular}
\end{center}

在强关联区域，HF 自洽性在数学上仍可收敛，但物理上已定性错误——多个行列式近简并，
没有单一行列式能够描述系统。此时 SQD 的威力恰恰体现在：它不对波函数形式做先验假设，
只要采样到的 bitstrings 包含了重要的组态，精确对角化就能恢复正确的基态能量。

\end{document}
